% PDF 138、140--143；原书第 108、110--113 页。
\subsection{自动设计模板——连续搜索}
构造提示的主要目标，是让大语言模型高效地完成任务。提示是否便于人理解，并非必备条件，因此无需把它限定在可供人理解的自然语言范围内（\refcite{ch03-li-2019}{Li et al., 2019）}。为此，研究者开始考察连续提示，又称软提示（soft prompt），直接在模型的嵌入空间中进行提示。

\begin{keypoint}
连续提示解决了离散提示的两项主要限制。
\begin{enumerate}
\item 模板中词的嵌入无需对应英语等自然语言中词的嵌入。
\item 模板不再局限于由预训练语言模型的参数来表示。模板可以拥有自己的参数，并使用下游任务的训练数据微调这些参数。
\end{enumerate}
\end{keypoint}

\subsubsection{前缀调优}
\refcite{ch03-li-liang}{Li 和 Liang（2021）}首次提出前缀调优（prefix tuning）。受提示的上下文学习效果启发（见第~\ref{sec:ch03-icl}~节），他们在输入文本开头添加了针对特定任务的“虚拟词元”（virtual token），见图~\ref{fig:ch03-prefix-tuning}。这些向量不代表实际词元。初始化时为它们设置适当的维数，使语言模型能像关注硬词元一样关注它们。随后，这些向量可作为连续变量参与训练，而硬词元的表示是固定的。语言模型由此可以借助前缀调优获知任务的性质，同时利用提示文本中明确的离散特征。

\refcite{ch03-li-liang}{Li 和 Liang（2021）}报告说，在数据量较少时，他们的前缀调优实验优于微调，与使用全量数据的微调相比也颇有竞争力。研究者在摘要任务中对 BART（\refcite{ch03-lewis}{Lewis et al., 2019）}做前缀调优，在表格到文本任务中对 GPT-2（\refcite{ch03-radford-2019}{Radford et al., 2019）}做同样的处理。与已有的适配器方法和全量数据微调相比，前缀调优在这两项任务上都表现良好。这些结果显示，前缀调优可用于不同类型的语言模型，实验具体展示了它在编码器—解码器模型和自回归模型上的适用性。

AutoPrompt 用训练数据集在离散空间中按梯度搜索，优化一组离散提示（见第~3.3.3.3~节）。与之类似，前缀调优用训练数据学习一组连续向量，即前缀，以最大化下式。
\begin{equation}
\max_{\phi}\log p_{\phi}(y\mid x)
  =\sum_{i\in Y_{\mathrm{idx}}}\log p_{\phi}(z_i\mid h_{<i}).
\tag{3.12}\label{eq:ch03-prefix-objective}
\end{equation}
通常，$p_{\phi}$ 由大语言模型可训练的参数 $\phi$ 决定。这里大语言模型的参数固定，因此用 $P_{\theta}$ 表示可训练的前缀参数 $\theta$。$h_i$ 是第 $i$ 个时间步各层激活值（包括前缀处激活值）的拼接。前缀调优利用从左到右生成的自回归语言模型。顾名思义，学得的向量放在模型序列最左侧，前缀的影响在模型的各个固定层中从左向右传播。

\sourcefigure{fig03_13.pdf}{上半图是预训练再微调。每做一项下游任务，都要复制并调优预训练模型的参数。有些语言模型规模很大，因此这样做开销显著，也有技术上的困难。下半图是前缀调优。它冻结模型参数，仅调优特定于任务的前缀。只需替换经过调优的前缀，同一个大语言模型便可用于多项下游任务。这样可将各任务的处理方式做成可替换的模块，并减少总体存储开销。}{fig:ch03-prefix-tuning}

\subsubsection{混合提示与用离散词元初始化的提示}
\refcite{ch03-li-liang}{Li 和 Liang（2021）}指出，前缀参数的初始化和对学习率的敏感性，可能使前缀调优的训练不稳定。他们在该项工作中将前缀参数化，以代替使用大型前馈神经网络生成的较小矩阵\footnote{译注：Li 和 Liang 论文第 4.3 节的做法是将较小的矩阵输入前馈网络来表示前缀，训练后仅保存得到的前缀。见 \url{https://aclanthology.org/2021.acl-long.353/}。}。另一种初始化连续词元的方法，是使用预先选定的离散词元。这些词元可以经由前文介绍的离散模板自动搜索学得（如 \refcite{ch03-zhong}{Zhong et al., 2021）}，也可以手工定义。在实体关系知识探查中，以这些词元为起点学习连续词元，已显示出良好前景（\refcite{ch03-qin}{Qin and Eisner, 2021}；\refcite{ch03-zhong}{Zhong et al., 2021）}。

\refcite{ch03-liu-2021}{Liu 等人（2021）}提出了另一种有前景的自动提示模板设计方法，用硬词元和软词元共同组成最终模板。图~\ref{fig:ch03-hybrid} 比较了离散提示模板（a）和这类混合提示（b）的生成过程。在 BERT 的知识探查任务中，他们的 P-tuning 方法生成的提示模板优于其他所有参与比较的离散提示模板。使用 ALBERT（\refcite{ch03-lan}{Lan et al., 2020）}时，这种方法在少样本 SuperGLUE 基准上达到了当时的最佳水平。

\sourcefigure{fig03_14.pdf}{对“The capital of Britain is [MASK]”搜索提示的例子。已知上下文（深灰色，“Britain”）和目标（白色方框，“[MASK]”），浅灰色区域表示提示词元。在（a）中，提示生成器只能收到离散奖励。在（b）中，可以采用不同方法优化伪提示和提示编码器。在某些情况下，加入少量与任务有关的锚定词元，如（b）中的“capital”，还可提高下游任务的表现。}{fig:ch03-hybrid}

\subsection{基于提示的微调}
\label{sec:ch03-prompt-finetuning}
前面介绍的提示方法，都固定用于推理的模型，通过改变提示或前缀来优化结果。另一种方法则解除对模型参数的冻结，将输入与输出对按固定模板组织成数据集，再用传统的反向传播方法微调模型参数。这种方法称为基于提示的微调（prompt-based fine-tuning）。仍以图~\ref{fig:ch03-formal-prompt} 的例子为例。
\begin{quote}
\texttt{Cannot watch this movie. This is [z].}
\end{quote}
除了可以用这个模板推理，还可以调优模型，使它准确预测人工指定的 \texttt{[z]} 值。要做好基于提示的微调，可能需要大量以同一模板组织的输入与输出对。不过，要获得良好表现，通常比常规 PTFT 所需的有标签样例少（参见第~\ref{sec:ch03-tutorial-one}~节的教程）。

这种方法既适用于完形填空式提示，也适用于前缀式提示。采用前者时，用相同形式的提示填入不同的输入与输出对来微调模型，使模型预测答案槽位中的词元。采用后者时，在答案的前缀上微调模型，使它逐步预测后续答案的下一个词元，答案可为任意长度。微调完成后，按标准的基于提示的学习流程用模型推理。

\begin{keypoint}
基于提示的微调是指令调优（instruction tuning）的基本技术，而指令调优是开发 ChatGPT 等先进聊天机器人的重要步骤。第~\ref{sec:ch04-instruction}~节将详细讨论指令调优。这里先指出，基于提示的训练一般优于零样本或少样本的基于提示的学习，代价是可能需要大量计算资源。它在许多情况下也优于 PTFT，尤其是在数据较少时。
\end{keypoint}
